% ---------------------------------------------------------------- GABARITO E RESUMO FINAL
\clearpage
\titulo{Gabarito e resumo}{Gabarito e resumo das 15 questões}
Use esta página para conferir suas respostas e para revisar a ideia que resolve cada questão.

{\small\renewcommand{\arraystretch}{1.26}
\begin{center}
\begin{tabular}{@{}c>{\raggedright\arraybackslash}p{2.9cm}>{\raggedright\arraybackslash}p{10.4cm}c@{}}
\toprule
\textbf{Questão} & \textbf{Área} & \textbf{A ideia que resolve} & \textbf{Resposta}\\
\midrule
91 & \bolinha{corB}\ Eletricidade & A força elétrica depende de $1/d^2$: distância menor, força muito maior. & \textbf{D}\\
92 & \bolinha{corD}\ Energia e calor & Somar as potências (9\,000\un{W}) e dividir por 0,3 para obter Btu/h. & \textbf{E}\\
94 & \bolinha{corA}\ Mecânica & Momento total zero: a partícula mais leve sai mais rápida e leva mais energia (4 partes de 5). & \textbf{E}\\
95 & \bolinha{corB}\ Eletricidade & Dois caminhos em 12\un{V}; em cada um, a tensão se divide entre os resistores ($U = R \cdot i$). & \textbf{D}\\
96 & \bolinha{corC}\ Ondas e óptica & Eco: a onda vai e volta (150\un{km}); tempo = distância $\div$ velocidade. & \textbf{E}\\
99 & \bolinha{corB}\ Eletricidade & A e B em paralelo, em série com C e D; lâmpada queimada interrompe o ramo. & \textbf{C}\\
104 & \bolinha{corD}\ Energia e calor & Abaixo do ponto triplo, o sólido passa direto a gasoso: sublimação. & \textbf{C}\\
111 & \bolinha{corD}\ Energia e calor & Urânio é limitado (não renovável) e a usina nuclear não depende do clima. & \textbf{A}\\
113 & \bolinha{corC}\ Ondas e óptica & Frequência e comprimento são inversamente proporcionais: $f_1 \cdot L_1 = f_2 \cdot L_2$. & \textbf{D}\\
114 & \bolinha{corD}\ Energia e calor & $1\un{kWh} = 3\,600\,000\un{J}$; dividir pela distância em metros e pelos 4 passageiros. & \textbf{B}\\
118 & \bolinha{corA}\ Mecânica & Na linha de separação, a pressão é igual dos dois lados: $d_A \cdot h_A = d_B \cdot h_B$. & \textbf{A}\\
121 & \bolinha{corC}\ Ondas e óptica & Intensidade = potência $\div$ área da esfera; a folha recebe só a sua área e absorve 60\%. & \textbf{E}\\
128 & \bolinha{corC}\ Ondas e óptica & As gotículas da neblina espalham a luz do farol para todos os lados. & \textbf{A}\\
130 & \bolinha{corA}\ Mecânica & Mesma vazão por uma saída menor: a velocidade da água aumenta. & \textbf{A}\\
133 & \bolinha{corB}\ Eletricidade & Energia = potência $\times$ tempo; comparar a energia em espera com a total. & \textbf{D}\\
\bottomrule
\end{tabular}
\end{center}}

\nota{Respostas obtidas pelas resoluções deste material. Confira com o gabarito oficial do Inep quando ele for divulgado.}

\subsecao{Cinco hábitos que resolveram esta prova}
\begin{enumerate}[label=\protect\selo{\arabic*},leftmargin=2em,itemsep=3pt]
  \item \textbf{Ler a pergunta até o fim.} ``Iguais e \emph{não nulas}'' (99), ``\emph{não renovável} que \emph{independe} do clima'' (111),
        ``energia \emph{em espera} comparada com o \emph{total}'' (133).
  \item \textbf{Conferir as unidades antes da conta.} W e Btu/h (92), kWh e joule (114), cm² e m² (121).
  \item \textbf{Perguntar ``o que muda e o que não muda?''} A distância muda, a carga não (91). A área da saída muda, a vazão não (130).
        A quantidade de líquido B não muda (118).
  \item \textbf{Usar o bom senso para conferir.} Corda mais curta tem de dar som mais agudo (113). A onda do radar é rapidíssima:
        o tempo tem de ser minúsculo (96).
  \item \textbf{Tirar o dado certo da figura.} O 60\% em 650\un{nm} (121), o 880\un{Hz} no quadro (113), o valor 90 comparado às barras (114).
\end{enumerate}
